\documentclass[10pt,a4paper]{article}
\usepackage[T1]{fontenc}
\usepackage{lmodern}
\usepackage[margin=22mm]{geometry}
\usepackage{amsmath,amssymb,booktabs,longtable,hyperref,fancyhdr,enumitem,graphicx}
\setlength{\parindent}{0pt}\setlength{\parskip}{.5em}
\setlength{\emergencystretch}{3em}
\pagestyle{fancy}\fancyhf{}
\fancyhead[L]{PCM / GS-PnP SSDU expansion}
\fancyhead[R]{8 November 2026}
\fancyfoot[C]{\thepage}\setlength{\headheight}{14pt}
\begin{document}
{\LARGE\bfseries Expanded GS-PnP SSDU tuning}\\[4pt]
{\large Dedicated preparation and experiment report}\\[4pt]
\textbf{Current status: Prepared, not run.}
\par\medskip\hrule
\section{Executive summary}
The previous ground-truth-free GS-PnP selection was
\(\lambda=0.1\), denoiser \(\sigma=0.075\) and
initial \(\tau_0=100/3\). Each was at the top of
the previous finite search range. The five-image held-out
measurement MSE was 0.0043654339.
We prepare 27 new-grid configurations, including the
old winning control and 26 previously untested combinations.
Selecting among them must use only archived validation
measurements, never image ground truth or the previously
inspected 450-image test scores.

\section{Tuning algorithm}
Use the same five NPL-noisy validation images:
CIGS 0733, Organic 0008, Perovskite 0283,
Reference 0768 and Silicone 0589.
For each image, retain the first 4,096 coarse rows
of the 13,107 acquired Hadamard measurements in
every reconstruction fold. Divide the remaining
9,011 acquired fine rows into the four previously
frozen disjoint SSDU holdouts.

For candidate \(\theta=(\lambda,\sigma,\tau_0)\),
reconstruct from fit rows \(\Theta_{ik}\)
and score held-out rows \(\Lambda_{ik}\):
\[
L(\theta)=\frac15\sum_{i=1}^{5}\frac14
\sum_{k=1}^{4}
\frac{\|A_{\Lambda_{ik}}\widehat{x}_{\theta,ik}
-y_{\Lambda_{ik}}\|_2^2}{|\Lambda_{ik}|}.
\]
Keep the original normalised PCM operator and summed
half-squared residual fidelity, without a fold-size
rescaling. Use the same pretrained GS-DRUNet and
paper-final \(z\) image.

\subsection*{Frozen parameter grid}
\begin{center}\begin{tabular}{ll}\toprule
Parameter & Candidate values\\\midrule
\(\lambda\) & 0.10, 0.20, 0.40\\
Denoiser \(\sigma\) & 0.075, 0.10, 0.15\\
Initial \(\tau_0\) & 33.3333, 50, 75\\\bottomrule
\end{tabular}\end{center}
Freeze 100 maximum accepted steps, 20 line-search attempts,
\(\gamma=0.1\), \(\eta=0.9\) and stopping
tolerance \(10^{-6}\). Tune \(\lambda\) and \(\tau_0\)
independently. Record accepted iterations,
rejections, final \(\tau\), effective
\(\lambda\tau\), and time per reconstruction.

\section{Step-by-step plan}
\begin{enumerate}[leftmargin=1.5em,itemsep=.2em]
\item Verify the frozen acquisition hashes, original four fold
masks, solver equivalence and data-fidelity scale.
Freeze source hashes and the candidate grid.
\item Reuse the original 20 fold scores of
\texttt{gs\_018} under control ID \texttt{gs\_exp\_001},
keeping original SHA256 provenance.
\item Score 26 new configurations across five images
and four folds: 520 new GS reconstructions.
Each can take up to 100 accepted iterations,
with additional backtracking work.
\item Record all held-out losses and solver diagnostics.
Exclude configurations with any failed fold.
If more than nine candidates fail, withhold selection.
\item Select the minimum balanced SSDU loss,
breaking ties within \(10^{-10}\) by grid order.
Freeze the candidate and every result hash.
\item Check whether the winning parameters are still
at the expanded upper bounds.
Only after freezing, consider a separately labelled
retrospective image-quality comparison.
\item Refresh this report with actual scores and
numerical diagnostics. Do not tune to test PSNR.
\end{enumerate}

\section{Experiment running}
The expanded sweep has not begun. The intended GPU is
tuni's RTX 4090, using an isolated CUDA UUID and
deterministic cuBLAS/cuDNN. The original float64
noisy measurements and mask indices are reused,
while neural solvers receive their required float32
values. The scoring code does not load targets.
Case-level JSON files allow safe resumption.
The old tuned-5 and tuned-50 selection records
and the 450-image test benchmark are preserved.

\section{Results and deliverables}
\textbf{Run status:} 1 valid, 0 invalid, 26 untested/incomplete candidates. The expanded parameter sweep has not been completed or frozen. The old control's held-out MSE was 0.0043654339. No higher-value configuration has yet been selected. This is a preparation report, not a result claiming improved GS-PnP.

\small
\begin{longtable}{@{}lrrrrl@{}}
\toprule
Candidate&\(\lambda\)&\(\sigma\)&\(\tau_0\)&
SSDU MSE&Status\\\midrule\endfirsthead
\toprule
Candidate&\(\lambda\)&\(\sigma\)&\(\tau_0\)&
SSDU MSE&Status\\\midrule\endhead
\texttt{gs\_exp\_001} & 0.10 & 0.075 & 33.3333 & 0.00436543 & reused control \\
\texttt{gs\_exp\_002} & 0.10 & 0.075 & 50.0000 & -- & not run \\
\texttt{gs\_exp\_003} & 0.10 & 0.075 & 75.0000 & -- & not run \\
\texttt{gs\_exp\_004} & 0.10 & 0.100 & 33.3333 & -- & not run \\
\texttt{gs\_exp\_005} & 0.10 & 0.100 & 50.0000 & -- & not run \\
\texttt{gs\_exp\_006} & 0.10 & 0.100 & 75.0000 & -- & not run \\
\texttt{gs\_exp\_007} & 0.10 & 0.150 & 33.3333 & -- & not run \\
\texttt{gs\_exp\_008} & 0.10 & 0.150 & 50.0000 & -- & not run \\
\texttt{gs\_exp\_009} & 0.10 & 0.150 & 75.0000 & -- & not run \\
\texttt{gs\_exp\_010} & 0.20 & 0.075 & 33.3333 & -- & not run \\
\texttt{gs\_exp\_011} & 0.20 & 0.075 & 50.0000 & -- & not run \\
\texttt{gs\_exp\_012} & 0.20 & 0.075 & 75.0000 & -- & not run \\
\texttt{gs\_exp\_013} & 0.20 & 0.100 & 33.3333 & -- & not run \\
\texttt{gs\_exp\_014} & 0.20 & 0.100 & 50.0000 & -- & not run \\
\texttt{gs\_exp\_015} & 0.20 & 0.100 & 75.0000 & -- & not run \\
\texttt{gs\_exp\_016} & 0.20 & 0.150 & 33.3333 & -- & not run \\
\texttt{gs\_exp\_017} & 0.20 & 0.150 & 50.0000 & -- & not run \\
\texttt{gs\_exp\_018} & 0.20 & 0.150 & 75.0000 & -- & not run \\
\texttt{gs\_exp\_019} & 0.40 & 0.075 & 33.3333 & -- & not run \\
\texttt{gs\_exp\_020} & 0.40 & 0.075 & 50.0000 & -- & not run \\
\texttt{gs\_exp\_021} & 0.40 & 0.075 & 75.0000 & -- & not run \\
\texttt{gs\_exp\_022} & 0.40 & 0.100 & 33.3333 & -- & not run \\
\texttt{gs\_exp\_023} & 0.40 & 0.100 & 50.0000 & -- & not run \\
\texttt{gs\_exp\_024} & 0.40 & 0.100 & 75.0000 & -- & not run \\
\texttt{gs\_exp\_025} & 0.40 & 0.150 & 33.3333 & -- & not run \\
\texttt{gs\_exp\_026} & 0.40 & 0.150 & 50.0000 & -- & not run \\
\texttt{gs\_exp\_027} & 0.40 & 0.150 & 75.0000 & -- & not run \\
\bottomrule

\end{longtable}
\normalsize
\IfFileExists{figures/candidate_losses.pdf}{\begin{figure}[ht]\centering\includegraphics[width=.95\linewidth]{figures/candidate_losses.pdf}\caption{Validation-only SSDU loss across completed candidates.}\end{figure}}{}
\IfFileExists{figures/parameter_slices.pdf}{\begin{figure}[ht]\centering\includegraphics[width=.95\linewidth]{figures/parameter_slices.pdf}\caption{Expanded lambda, sigma and initial step-size slices of the validation-only risk.}\end{figure}}{}
The dedicated folder holds the original grid, preflight,
resumable numerical code, case-level results, frozen selection
and an editable report in both TeX and PDF.

\section{Conclusion}
An expanded parameter grid can establish whether the
original SSDU measurement-risk search was too narrow.
It cannot by itself establish globally optimal reconstruction,
or guarantee preservation of thin structures.
Higher parameter values may worsen numerical stability.
Correlated PCM noise and unmeasured Hadamard directions
also limit the link between SSDU risk and image quality.

\paragraph{Primary papers.}
Yaman et al., \emph{Self-Supervised Learning of
Physics-Guided Reconstruction Neural Networks without
Fully-Sampled Reference Data}, Magnetic Resonance in Medicine
(2020), acquired measurement partitioning:
\url{https://doi.org/10.1002/mrm.28378}.

Hurault, Leclaire and Papadakis,
\emph{Gradient Step Denoiser for Convergent Plug-and-Play},
ICLR (2022), gradient-step update and backtracking:
\url{https://openreview.net/pdf?id=fPhKeld3Okz}.
\end{document}